\chapter{SES / UF-Core Security and Protected Structural Boundary}
\section{Security Objectives}
\meta{State the mandatory security objectives inherited into TFE.}{UF-SP, SES Summary, deployment environment, domain parameters, threat model.}{Normative security goals for TFE.}

TFE shall inherit and enforce the following normative security objectives:
\begin{enumerate}[leftmargin=1.2cm]
    \item protect the UF Kernel implementation against practical reverse-engineering and tampering;
    \item protect UF-SP and derivative protected structural representations from disclosure or undetected modification;
    \item support multi-tenant, multi-environment deployment under a single deterministic model; and
    \item implement mandatory threat responses when specific security signals are observed.
\end{enumerate}

\section{UF-SP, SES Summary, and Redaction Policy}
\meta{Define what is protected and what may be plaintext.}{Kernel outputs, diagnostics, redaction policy.}{UF-SP serialization rule and SES Summary whitelist.}

UF-SP shall include, at minimum:
\begin{itemize}
    \item SEV sequences from L0;
    \item gate boundaries and attributes from L1;
    \item interpretive vectors from L2;
    \item resonance curves from L3;
    \item DSF vectors from L4; and
    \item validation/hardening diagnostics.
\end{itemize}

UF-SP must be serialized deterministically prior to encryption.

Plaintext outputs outside the protected boundary shall be restricted to the SES Summary whitelist:
\begin{itemize}
    \item $S(UF)$;
    \item $R(UF)$;
    \item GSS;
    \item DSS;
    \item DSF-Stability scalar;
    \item collapse flags;
    \item SafeMode indicator.
\end{itemize}

No raw SEV, gate internals, full DSF vectors, per-gate diagnostics, resonance curves, or protected research event tapes may cross the boundary in plaintext.

\section{Domain Parameters}
\meta{Define the canonical domain-binding tuple for TFE.}{Asset identity, window identity, versions, tenant/environment identifiers.}{DomainParams tuple and equality rule.}

For TFE, the canonical domain-binding tuple is
\[
\operatorname{DomainParams} = (symbol, window\_start, window\_end, timeframe, uf\_version, ses\_version, tenant\_id, engine\_variant, environment).
\]

A TFE implementation may add a non-sensitive adapter/profile identifier, but it shall not remove any required field. Two evaluations are domain-equal if and only if all tuple members match exactly.

\section{Key Hierarchy and Envelope Format}
\meta{Define the mandatory key derivation and envelope structure.}{$K_{root}$, DomainParams, UF-SP, SES Summary.}{$K_{domain}$ and a valid SES Envelope.}

TFE shall use a two-layer key hierarchy:
\begin{align*}
K_{domain} &= HKDF(K_{root},\operatorname{DomainParams}).
\end{align*}

The SES Envelope shall be structured as
\[
\operatorname{Envelope} = \{header, summary, ciphertext\}
\]
with
\begin{align*}
header &= \{\operatorname{DomainParams}\},\\
summary &= \{S,R,GSS,DSS,DSFStability,collapse\_flags,SafeMode\},\\
ciphertext &= C.
\end{align*}

\section{Encode / Decode Algorithms}
\meta{Specify the mandatory envelope transformation semantics.}{UF-SP, DomainParams, $K_{root}$, deterministic serializer.}{SES Envelope or verified UF-SP.}

Let $\operatorname{Serialize}(\cdot)$ and $\operatorname{Deserialize}(\cdot)$ be canonical deterministic serializers.

\textbf{Encode:}
\begin{align*}
header &\leftarrow \operatorname{DomainParams},\\
summary &\leftarrow \operatorname{Redact}(UF\mbox{-}SP),\\
K &\leftarrow HKDF(K_{root},\operatorname{DomainParams}),\\
A &\leftarrow \operatorname{Serialize}(header,summary),\\
P &\leftarrow \operatorname{Serialize}(UF\mbox{-}SP),\\
C &\leftarrow SIV\mbox{-}Encrypt(K,P,A),\\
\operatorname{Envelope} &\leftarrow \{header,summary,C\}.
\end{align*}

\textbf{Decode:}
\begin{align*}
K &\leftarrow HKDF(K_{root},\operatorname{DomainParams}),\\
A &\leftarrow \operatorname{Serialize}(header,summary),\\
P &\leftarrow SIV\mbox{-}Decrypt(K,C,A),\\
UF\mbox{-}SP &\leftarrow \operatorname{Deserialize}(P).
\end{align*}

Decryption shall fail if any element of DomainParams differs between encode and decode.

\section{Protected Research Profile (PRP) Rule}
\meta{Define how richer structural representations may be used inside TFE without violating SES.}{SES Envelopes, PRP authorization, controlled research execution.}{Admissible protected use of richer structural payloads.}

In PRP, L5 may derive richer representations --- including structural event tapes, mode-amplitude histories, or protected state checkpoints --- from decrypted UF-SP \emph{only} inside a controlled boundary satisfying all of the following:
\begin{enumerate}[leftmargin=1.2cm]
    \item PRP mode is explicitly designated and logged;
    \item decryption occurs only after successful domain and environment verification;
    \item derived protected representations are treated as protected structural payloads and must themselves be enveloped or remain memory-resident only;
    \item plaintext protected representations are never written to disk, logs, or network outside the protected boundary;
    \item ordinary TFE consumers still receive only SES Summary fields or explicitly approved coarse decisions.
\end{enumerate}

\section{Threat and Response Matrix}
\meta{Define mandatory behavior when specific security signals are observed.}{Security signals from kernel, SES, runtime, and hosting environment.}{Local and system responses.}

\renewcommand{\arraystretch}{1.15}
\begin{longtable}{p{0.10\textwidth}p{0.26\textwidth}p{0.26\textwidth}p{0.28\textwidth}}
\toprule
Threat ID & Description & Detection Signals & Required Response \\
\midrule
T1 & Code extraction attempt & Debug / introspection flags, breakpoint patterns, known analysis environments where available & Local: set SafeMode, restrict outputs to coarse summaries, disable UF-SP export. System: log security event and notify operator. \\
T2 & Code tampering & Bundle / manifest hash mismatch & Local: refuse normal operation, enter compromised state. System: quarantine/replace instance and route workloads away. \\
T3 & UF-SP exfiltration & Request for raw UF-SP without SES envelope or invalid envelope fields & Local: deny request. System: log and optionally throttle/block caller. \\
T4 & Envelope replay / misbinding & DomainParams mismatch or SIV authentication failure & Local: reject decryption and return no UF-SP. System: log replay suspicion and optionally rate-limit/review credentials. \\
T5 & Engine rollback / clone & UF version or engine variant inconsistent with expected deployment & Local: flag anomaly, optionally degrade outputs. System: prevent registration without explicit approval. \\
T6 & Adversarial environment & Missing integrity guarantees, unknown host, missing KMS access where required & Local: minimum safe mode only. System: deny node registration for production workloads. \\
T7 & Rogue tenant / interface abuse & Excessive or abnormal query patterns per domain & Local: enforce rate limits. System: throttle or suspend tenant identity. \\
T8 & Domain parameter manipulation & Repeated attempts to alter symbol/timeframe/tenant/environment for cross-domain reuse & Local: treat as error, never decrypt under mismatch. System: flag suspicious patterns. \\
T9 & Data integrity tampering & Non-monotone timestamps, impossible bars, structural invalidity & Local: reject input, emit no normal envelope. System: log anomaly and request upstream verification. \\
T10 & Collapse-induction denial of service & Repeated collapse-driving or SafeMode-triggering inputs & Local: limit repeated processing. System: stronger throttling, source isolation, or blacklist as approved. \\
\bottomrule
\end{longtable}

\section{Deployment Profiles Applicable to TFE}
\meta{State the deployment modes under which TFE may operate.}{Hosting environment and protection capabilities.}{Profile obligations.}

TFE shall support the following profiles:
\begin{itemize}
    \item \textbf{Profile C (Cloud Service):} controlled cloud compute, protected storage, KMS/HSM access, SES at rest and in transit.
    \item \textbf{Profile E (Edge / On-Prem Engine):} same envelope and logging rules, optional TPM/HSM support, restricted local admin interfaces.
    \item \textbf{Profile D (Constrained Device Wrapper):} device sends raw series to protected engine and receives only summary or coarse decisions.
\end{itemize}

Optional informative profiles such as shadow approximators or split-brain deployments may be implemented only if their semantics, security posture, and validation obligations are separately specified.

\section{Compliance Requirements}
\meta{State the minimum conditions for claiming conformance to the TFE security boundary.}{Kernel boundary, envelope controls, plaintext restrictions, threat responses, logging.}{Conformance rule.}

A TFE implementation claiming conformance to this specification shall:
\begin{enumerate}[leftmargin=1.2cm]
    \item implement UF L0--L4 inside a clearly defined kernel boundary;
    \item implement SES/UF-Core envelope protection over UF-SP and equivalent protection over any PRP-derived protected representations;
    \item ensure that all protected structural payloads at rest and in transit are encrypted and authenticated via SES Envelopes or equivalent profile-approved wrappers;
    \item restrict plaintext outputs to SES Summary fields only outside the protected boundary;
    \item implement all mandatory threat responses T1--T10; and
    \item provide a chain-of-custody log containing versions, domain parameters, environment, and all security events without exposing UF-SP.
\end{enumerate}
